\subsection{Вещественные числа: обзор IEEE 754}
\label{lecture02-float}
Вещественное число записывают в виде $x=(-1)^s M\,2^E$.
Для ненулевого нормализованного двоичного числа $1\le M<2$.
Например:
\[12.625_{10}=1100.101_2=1.100101_2\cdot2^3.\]
Точка «плавает», потому что порядок $E$ меняет масштаб.
Конечное число битов позволяет хранить только конечный набор значений.

\begin{center}
\begin{tabular}{|l|c|c|c|c|}
\hline
Формат & Знак $s$ & Порядок $e$ & Дробь $f$ & Смещение \\
\hline
\texttt{float32} (binary32) & 1 & 8 & 23 & 127 \\
\texttt{double} (binary64) & 1 & 11 & 52 & 1023 \\
\hline
\end{tabular}
\end{center}
Размеры полей даны в битах. В распространённых реализациях C
\texttt{float} и \texttt{double} используют эти форматы, но это
свойство реализации: один лишь язык C не гарантирует IEEE 754.

Для нормализованных чисел поле $f$ трактуется как дробь:
\[0.f_2=\sum_{i=1}^{q} f_i2^{-i},\]
где $q$ --- число битов поля. Значение числа:
\[
x=(-1)^s(1+0.f_2)2^{e-\mathrm{bias}}.
\]
Ведущая единица не хранится. Поэтому точность significand равна 24
значащим двоичным разрядам у binary32 и 53 у binary64.
Не следует путать всю значащую часть $1.f_2$ с хранимым полем $f$.

\subsubsection*{Нормальные, субнормальные и особые значения}
Для binary32:
\begin{itemize}
    \item $1\le e\le254$: нормализованные числа, $E=e-127$.
    \item $e=0$, $f=0$: $+0$ или $-0$ в зависимости от $s$.
    \item $e=0$, $f\ne0$: субнормальные (денормализованные) числа,
    без скрытой единицы:
    \[x=(-1)^s(0.f_2)2^{-126}.\]
    \item $e=255$, $f=0$: $+\infty$ или $-\infty$.
    \item $e=255$, $f\ne0$: NaN; кодирует результат, не являющийся числом.
\end{itemize}
Наименьшее положительное нормальное binary32 равно $2^{-126}$,
субнормальное --- $2^{-149}$. Субнормальные числа обеспечивают
постепенный переход к нулю.
Для binary64 нормальное поле порядка лежит в $1\ldots2046$;
$e=0$ выделено для нулей и субнормальных, $e=2047$ --- для бесконечностей
и NaN. Истинный порядок субнормальных равен $-1022$.

\subsubsection*{Пошаговое кодирование $-12.625$}
\begin{enumerate}
    \item Знак отрицательный: $s=1$.
    \item Модуль: $1100.101_2=1.100101_2\cdot2^3$.
    \item Истинный порядок $E=3$, хранимый $e=3+127=130=10000010_2$.
    \item Поле дроби: \texttt{10010100000000000000000} (23 бита).
    \item Результат: \texttt{1 10000010 10010100000000000000000},
    или \texttt{0xC14A0000}.
\end{enumerate}
Проверка: $-(1+2^{-1}+2^{-4}+2^{-6})\cdot2^3=-12.625$.
Это точное представление. Hex-запись описывает 32-битное слово;
порядок его байтов в памяти --- отдельный вопрос. В little-endian
памяти байты идут как \texttt{00 00 4A C1}.

\subsubsection*{Машинная точность и округление}
Машинный эпсилон здесь определяется как расстояние от 1 до следующего
представимого числа: $\varepsilon=2^{-23}$ для binary32,
$2^{-52}$ для binary64. Это \emph{не} наименьшее положительное число
и не универсальная абсолютная погрешность.
В интервале $[2^E,2^{E+1})$ шаг нормальных положительных binary32
равен $2^{E-23}$: с ростом масштаба увеличивается абсолютный шаг.

Конечная двоичная дробь после сокращения имеет знаменатель --- степень
двойки. Поэтому $0.5=1/2$ хранится точно, а $0.1=1/10$ --- приближённо.
Обычный режим округления IEEE 754 --- к ближайшему, при равенстве
расстояний к варианту с чётным младшим сохраняемым битом
(round to nearest, ties to even).

\textit{Учебный пример:} сохраняем 2 бита после двоичной точки.
Число $1.011_2=1.375$ находится ровно между $1.01_2=1.25$ и
$1.10_2=1.5$. Выбираем $1.10_2$: его младший бит равен 0.
Число $1.0101_2=1.3125$ ближе к $1.01_2$, поэтому округляется вниз.

\subsubsection*{Что важно для программирования}
\begin{itemize}
    \item \textbf{Поглощение:} в binary32 $2^{24}+1$ округляется к $2^{24}$.
    Ближайшие значения сверху различаются на 2.
    \item \textbf{Неассоциативность:} при округлении после каждой операции
    для $a=2^{24}$, $b=1$, $c=-2^{24}$ получаем
    $(a+b)+c=0$, но $a+(b+c)=1$.
    \item \textbf{Вычитание близких приближений} может усилить относительную
    ошибку исходных данных. Само вычитание близких точно представимых
    чисел не обязательно вносит ошибку.
    \item \textbf{Сравнение:} если по смыслу задачи достаточно близости,
    используют абсолютный и относительный допуски. Для конечных значений:
    \[|a-b|\le\max(\delta_{\rm abs},\delta_{\rm rel}\max(|a|,|b|)).\]
    Допуски выбирают по задаче, а не автоматически
    равными машинному эпсилону. Проверка точного равенства иногда уместна.
\end{itemize}
Алгоритмы арифметики, нормализация после вычитания, округление и подробные
примеры рассматриваются в лекции №3 (с.~\pageref{lecture03}).
В презентации лекции №2 они доступны как дополнительный блок.

\subsection{Кодирование машинных команд}
Программа тоже хранится в виде байтов. Процессор читает команду по адресу
счётчика команд, декодирует её, получает операнды и выполняет действие.
Следующий адрес обычно определяется длиной текущей инструкции;
переходы меняют обычную последовательность исполнения.

\textbf{Система команд} (ISA) определяет доступные операции, регистры,
форматы и способы адресации. Код операции (opcode) сообщает, что делать.
Другие поля могут задавать регистры, непосредственное значение, адрес
или смещение. Некоторые операнды и место результата бывают неявными.
Поэтому команда не обязана явно содержать три адреса и адрес следующей команды.

\textit{Учебная 16-битная команда} имеет формат
\[
\underbrace{\texttt{0011}}_{\text{opcode: 4}}
\ \underbrace{\texttt{010}}_{\text{dst: 3}}
\ \underbrace{\texttt{101}}_{\text{src: 3}}
\ \underbrace{\texttt{000000}}_{\text{резерв: 6}}.
\]
Зададим сами: opcode \texttt{0011} означает сложение,
операция $R_{dst}\gets R_{dst}+R_{src}$.
Тогда слово означает $R_2\gets R_2+R_5$ и имеет hex-код
\texttt{0x3540}. В этой модели доступны 16 кодов операций и 8 регистров.
Это учебный формат, а не кодировка x86.

В x86 инструкции имеют переменную длину (от 1 до 15 байт).
Например, в 64-битном режиме NASM-команда \texttt{mov eax, 1}
может кодироваться как \texttt{B8 01 00 00 00}:
opcode с номером регистра и 32-битное непосредственное значение.
\texttt{add eax, 2} имеет короткую кодировку \texttt{83 C0 02}.
Одинаковые исходные действия могут иметь несколько допустимых кодировок.
Длину и смысл байтов определяют ISA и режим исполнения.

\subsubsection*{Маски как инструмент обработки полей}
Пусть биты нумеруются справа с нуля и $0\le k<8$.
Для 8-битного слова маска $m=2^k$ выделяет бит $k$:
\begin{itemize}
    \item проверка: $(x\mathbin{\&}m)\ne0$;
    \item установка: $x\mathbin{|}m$;
    \item сброс: $x\mathbin{\&}(255\mathbin{\mathrm{XOR}}m)$;
    \item переключение: $x\mathbin{\mathrm{XOR}}m$.
\end{itemize}
Для $x=\texttt{0x5A}=\texttt{01011010}$ бит 2 равен 0;
его установка даёт \texttt{0x5E}, сброс бита 4 --- \texttt{0x4A}.
Поле из трёх битов, начинающееся с бита 2, извлекается как
\texttt{(x >> 2) \& 7} и равно 6.

\subsection*{Итоги и вопросы для самопроверки}
\begin{enumerate}
    \item Почему байт \texttt{FF} можно интерпретировать как 255 и как $-1$?
    \item Какие предположения нужны для расчёта объёма изображения или PCM?
    \item Почему два символа не обязательно занимают два байта?
    \item Чем перенос CF отличается от знакового переполнения OF?
    \item Почему $0.1$ не представляется конечной двоичной дробью?
    \item Что нужно знать, чтобы декодировать байты как машинную команду?
\end{enumerate}
